\documentclass[10pt,a4paper]{amsart}
\usepackage[margin=19mm]{geometry}
\usepackage{amsmath,amssymb,mathtools}
\usepackage[scr=rsfs,cal=euler]{mathalfa}
\usepackage{xfrac}
\usepackage[unicode,hidelinks,bookmarks=false]{hyperref}
\newcommand{\ie}{i.e.\ }
\newcommand{\cf}{cf.\ }
\newcommand{\resp}{respectively\ }
\newcommand{\Subsection}{\subsection}
\providecommand{\nobreakdash}{\nobreak\hbox{-}}
\setlength{\emergencystretch}{2em}
\multlinegap0pt
\usepackage{amssymb}
\usepackage{algorithm}\usepackage{algpseudocode}
\newtheorem{lemma}{Lemma}
\newcommand{\Mle}{\mathcal{M}_{\le r}}
\newcommand{\R}{\mathbb{R}}
\newcommand{\eps}{\varepsilon}
\newcommand{\range}{\operatorname{range}}
\newcommand{\rank}{\operatorname{rank}}
\title{Robust High-Order Projector-Splitting Integrators}
\author{Shiheng Zhang, Jingwei Hu}
\date{Source: arXiv:2608.17157v2 (CC BY 4.0)}
\begin{document}
\maketitle

\noindent\emph{Source and licence.} Robust High-Order Projector-Splitting Integrators. Version arXiv:2608.17157v2 (CC BY 4.0), distributed by arXiv under the Creative Commons Attribution 4.0 International licence, http://creativecommons.org/licenses/by/4.0/, as recorded in the arXiv licence metadata for this exact version and shown on the abstract page https://arxiv.org/abs/2608.17157v2. Full licence notice: \texttt{licenses/arxiv-2608.17157-CC-BY-4.0.txt}.\par\smallskip

\noindent\textit{Editorial scope.} Counted unit: the article's lemma lem:bounded-generators (source0.tex lines 1283--1298). The article's complete original proof of it is delivered with it (source0.tex lines 1300--1398, one \texttt{proof} environment of 99 lines). No mathematics is written, repaired or re-derived: the statement and the proof are the article's own lines, in the article's own notation, and no retained line is substituted or re-flowed.  Retained uncounted as the source-local context the proof needs: lines 436--440; lines 1236--1248.  Nothing was omitted from any retained span, so the package carries no omission record.  No retained line contains a citation, so the wrapper carries no bibliography and no bibliography entry.  No retained line contains a figure environment, an image-inclusion command or drawing code, so no image is delivered and the package declares no asset dependency.  Editorial formatting only, in addition: an AMS \texttt{amsart} preamble that restates the article's own declarations and macros used by the retained lines, namely the environments \texttt{lemma} on separate counters, together with the standard packages those lines need.  The retained environments print in this document as Lemma 1 in order of appearance, whereas the article prints the same items with the numbering of its own full document.  The complete native source of the article as distributed in the arXiv e-print for this exact version is bundled unchanged as \texttt{sources/207875/source0.tex}.
\begin{equation}\label{eq:assumption-BL}
 \|F(t,X)\|\le B,
 \qquad
 \|F(t,X)-F(t,Z)\|\le L_F\|X-Z\|.
\end{equation}

\begin{lemma}
\label{lem:lower-stratum-defect}
Let \(X\in\Mle\).  If \(x\) and \(y\) are unit vectors satisfying
\begin{equation}\label{eq:lower-normal-directions}
 x\perp\range(X),
 \qquad
 y\perp\range(X^\top),
\end{equation}
then
\begin{equation}\label{eq:lower-stratum-defect}
 |x^\top F(t,X)y|\le\eps_r.
\end{equation}
\end{lemma}

\begin{lemma}
\label{lem:bounded-generators}
For every \(t\in[0,T]\) and \(X\in\Mle\), there are matrices
\(G_L(t,X)\in\R^{m\times m}\), \(G_R(t,X)\in\R^{n\times n}\), and
\(E(t,X)\in\R^{m\times n}\) such that
\begin{equation}\label{eq:generator-decomposition}
 F(t,X)=G_L(t,X)X+XG_R(t,X)+E(t,X),
\end{equation}
with
\begin{equation}\label{eq:generator-bounds}
 \|G_L(t,X)\|_2,\ \|G_R(t,X)\|_2
 \le K_G:=(r+\sqrt r)L_F,
 \qquad
 \|E(t,X)\|\le E_G:=\sqrt{mn}\,\eps_r.
\end{equation}
\end{lemma}

\begin{proof}
Let \(q:=\rank(X)\).  If \(q=0\), then \(X=0\).  Lemma
\ref{lem:lower-stratum-defect}, applied to arbitrary unit vectors \(x,y\),
shows that every entry of \(F(t,0)\) in orthonormal coordinates is bounded
by \(\eps_r\).  Thus
\begin{equation}\label{eq:zero-generator-case}
 G_L(t,0):=0,\qquad G_R(t,0):=0,\qquad E(t,0):=F(t,0)
\end{equation}
satisfies \eqref{eq:generator-decomposition}--\eqref{eq:generator-bounds}.

Suppose \(q\ge1\), and take a compact singular value decomposition (SVD)
\begin{equation}\label{eq:generator-svd}
 X=U\Sigma V^\top
  =\sum_{i=1}^q\sigma_i u_i v_i^\top,
 \qquad \sigma_i>0.
\end{equation}
Define
\begin{align}
 X_i&:=X-\sigma_i u_i v_i^\top,                       \label{eq:delete-one}\\
 X_{ij}&:=X-\sigma_i u_i v_i^\top-\sigma_j u_j v_j^\top
 \quad(i\ne j).                                      \label{eq:delete-two}
\end{align}
Define
\begin{align}
 g_i^L&:=\frac{(I-UU^\top)(F(t,X)-F(t,X_i))v_i}{\sigma_i},
                                                               \label{eq:left-generator-column}\\
 g_i^R&:=\frac{(I-VV^\top)(F(t,X)-F(t,X_i))^\top u_i}{\sigma_i}.
                                                               \label{eq:right-generator-column}
\end{align}
The Lipschitz bound in \eqref{eq:assumption-BL} gives
\begin{equation}\label{eq:normal-generator-bounds}
 \|g_i^L\|_2,\ \|g_i^R\|_2\le L_F.
\end{equation}
With
\(G_\perp^L:=[g_1^L,\ldots,g_q^L]\) and
\(G_\perp^R:=[g_1^R,\ldots,g_q^R]\), it follows that
\begin{equation}\label{eq:normal-generator-matrix-bounds}
 \|G_\perp^L\|_2,\ \|G_\perp^R\|_2\le\sqrt q\,L_F.
\end{equation}

Define \(\Gamma\in\R^{q\times q}\) by
\begin{align}
 \Gamma_{ii}
 &:=\frac{u_i^\top(F(t,X)-F(t,X_i))v_i}{2\sigma_i},
                                                               \label{eq:gamma-diagonal}\\
 \Gamma_{ij}
 &:=\frac{u_i^\top(F(t,X)-F(t,X_{ij}))v_j}
          {\sigma_i+\sigma_j}
 \quad(i\ne j).                                      \label{eq:gamma-offdiagonal}
\end{align}
Since
\(\|X-X_{ij}\|=(\sigma_i^2+\sigma_j^2)^{1/2}
\le\sigma_i+\sigma_j\),
\begin{equation}\label{eq:gamma-bound}
 |\Gamma_{ii}|\le\frac{L_F}{2},\qquad
 |\Gamma_{ij}|\le L_F,\qquad
 \|\Gamma\|_2\le\|\Gamma\|\le qL_F.
\end{equation}
Now set
\begin{equation}\label{eq:generators-defined}
 G_L:=G_\perp^LU^\top+U\Gamma U^\top,
 \qquad
 G_R:=V(G_\perp^R)^\top+V\Gamma V^\top.
\end{equation}
Equations \eqref{eq:normal-generator-matrix-bounds} and
\eqref{eq:gamma-bound} yield
\begin{equation}\label{eq:generators-norm}
 \|G_L\|_2,\ \|G_R\|_2
 \le(\sqrt q+q)L_F\le K_G.
\end{equation}

It remains to bound
\begin{equation}\label{eq:generator-residual-defined}
 E:=F(t,X)-G_LX-XG_R.
\end{equation}
Complete \(U,V\) to orthogonal matrices \([U,U_\perp]\) and
\([V,V_\perp]\).  Direct substitution of
\eqref{eq:left-generator-column}--\eqref{eq:generators-defined} gives
\begin{align}
 (I-UU^\top)EVV^\top
 &=\sum_{i=1}^q (I-UU^\top)F(t,X_i)v_i v_i^\top,       \label{eq:residual-left-normal}\\
 UU^\top E(I-VV^\top)
 &=\sum_{i=1}^q u_i u_i^\top F(t,X_i)(I-VV^\top),      \label{eq:residual-right-normal}\\
 (U^\top EV)_{ii}
 &=u_i^\top F(t,X_i)v_i,                              \label{eq:residual-core-diagonal}\\
 (U^\top EV)_{ij}
 &=u_i^\top F(t,X_{ij})v_j\quad(i\ne j),              \label{eq:residual-core-offdiagonal}\\
 (I-UU^\top)E(I-VV^\top)
 &=(I-UU^\top)F(t,X)(I-VV^\top).                       \label{eq:residual-normal-normal}
\end{align}
For every scalar entry on the right-hand sides, the left and right vectors
are normal directions to \(X_i\), \(X_{ij}\), or \(X\), respectively.
Lemma~\ref{lem:lower-stratum-defect} therefore bounds every entry of
\([U,U_\perp]^\top E[V,V_\perp]\) by \(\eps_r\).  Hence
\begin{equation}\label{eq:generator-residual-bound}
 \|E\|\le\sqrt{mn}\,\eps_r=E_G,
\end{equation}
which completes the proof.
\end{proof}

\end{document}
